% Lösungen Einheit 4 von 4 – Lagebefunde im Sachzusammenhang (zusammenbau v0.1)
\einheitenkopf{Einheit 4 von 4 · Lagebefunde im Sachzusammenhang}

\erg{16}{a) nein – ein Bereich ist noch zu prüfen: die Koordinaten gegen die Wandmaße. \quad b) Lichtgerade $\vec{x} = (2 | 0 | 4) + t \cdot (1 | 2 | -1)$; $2t = 6$ gibt $t = 3$, Schattenpunkt $(5 | 6 | 1)$; $0 \le 5 \le 6$ und $0 \le 1 \le 3$ – der Schatten liegt auf der Wand. \quad c) Schritt I bestimmt den Schattenpunkt $(3 | 0 | 2)$ der Ecke $(4 | 2 | 3)$ als Durchstoßpunkt der Lichtgeraden mit der Wandebene; Schritt II zeigt, dass er innerhalb der Wandmaße liegt – der Schatten der Ecke fällt auf die Hauswand. \quad d) $4 + 4t = 6$ gibt $t = 0,5$; die Höhe dort ist $3 - 0,5 = 2,5$ m $> 2$ m – der Ball fliegt über das Netz und berührt es nicht. \quad e) $-5t^2 + 3t + 2 = 0$ gibt $t = 1$ (die Lösung $t = -0,4$ liegt vor dem Schlag); Auftreffpunkt $(4 | 11 | 0)$; $0 \le 4 \le 9$ und $0 \le 11 \le 18$ – der Ball trifft im Feld auf.}
\erg{17}{a) Die Gerade durch $T$ mit dem Richtungsvektor $\vec{v}$ mit der Dachebene $D$ schneiden; der Schatten reicht vom Fußpunkt $F$ bis zu diesem Schnittpunkt $Q$, seine Länge ist der Abstand $|\vec{FQ}|$ – nicht der Schnitt mit dem Boden.}
\erg{18}{a) Paula hat die Wandmaße nicht geprüft: $x = 2 > 1,8$ und $z = 3 > 2,5$ – der Schatten liegt nicht auf der Wand. \quad b) Licht läuft geradlinig; der Strahl, der an $P$ vorbeigestreift wäre, trifft dort auf die Wand, wo die Lichtgerade durch $P$ die Wandebene schneidet – dort fehlt das Licht.}
